\documentclass[11pt]{article}

% Change "review" to "final" for the camera-ready, or "preprint" for a
% non-anonymous version with page numbers.
\usepackage[review]{acl}

\usepackage{times}
\usepackage{latexsym}
\usepackage[T1]{fontenc}
\usepackage[utf8]{inputenc}
\usepackage{microtype}
\usepackage{inconsolata}
\usepackage{graphicx}
\usepackage{booktabs}
\usepackage{amsmath}
\usepackage{amssymb}
\usepackage{multirow}
\usepackage{xcolor}

\graphicspath{{figures/}}

% \todo{...} marks an unverified or untraced number/claim for the critics.
\newcommand{\todo}[1]{\textcolor{red}{\textbf{[TODO: #1]}}}

% ============================================================================
% This file is maintained by the `paper-writer` subagent across loop rounds.
% Industrial-track framing: a practitioner must choose ONE multilingual
% embedding model to deploy for chemistry-patent retrieval, and must measure
% cross-lingual robustness honestly (average recall hides cross-lingual
% collapse). The paper contributes (a) two patent-grounded multilingual/CLIR
% benchmarks, (b) new robustness metrics (CLIR-MRS / MRS and the cross-lingual
% RBO / confusion / separability family), and (c) a model-by-model analysis
% yielding a concrete deployment recommendation.
% NOTE: Every number in this paper MUST trace to a file under reports/.
% ============================================================================

\title{Measuring Cross-Lingual Robustness for Chemistry-Patent Retrieval:
Patent-Grounded Multilingual Benchmarks, New Metrics, and a Deployment Decision}

\author{Anonymous Submission \\
  EMNLP Industry Track \\
  \texttt{anonymous@example.com}}

\begin{document}
\maketitle

\begin{abstract}
A chemistry-patent search team that must deploy a single multilingual
embedding model typically ranks candidates by average Recall@10---the one
number that hides the failure they will ship. Average recall is inflated by
\emph{same-language} hits; the moment a query must reach a patent in another
language (the normal case in a patent family), recall collapses, and no two
language versions of the same question return the same documents. We make this
collapse measurable. We release two \emph{content-controlled},
patent-grounded multilingual chemistry-retrieval benchmarks built only from
human-translated patent text---an alias-graph benchmark (132 questions, 24
ChEBI compounds, with chemically-confusable hard negatives) and a cross-lingual
(CLIR) benchmark (137 questions over a 23{,}487-document shared corpus, with
Spanish as a pure ``no-home'' query language)---and a robustness-metric family
that reports cross-lingual recall, ranking consistency, confusion, and
embedding separability \emph{next to} recall. Across nine models, best
cross-lingual Recall@10 is only $0.50$, the same-language home advantage
reaches $+0.55$, five language versions of a query agree on barely a third of
their top-10 (cross-lingual RBO ceiling $0.39$ on the alias-graph benchmark,
$0.19$ on the cross-lingual benchmark), and a wrong but confusable compound
out-ranks the right one on $14$--$78\%$ of queries. We further quantify
\emph{what cross-linguality costs}: to find a foreign twin at the median you
scan $3.5\times$ as many documents as for a same-language copy
(\emph{cross-lingual reading-cost multiplier}, XRC50), a top-100 re-ranker
can recover at most $74\%$ of foreign twins (\emph{re-ranker recoverability
ceiling}, RRC@100), and $5.8\%$ are structurally unrecoverable by any re-ranker
($L_\infty$)---a measured share recoverable only by alignment, not re-ranking. We
show the cause is an embedding-level separability
deficit ($r(\text{cross-language AUC},\text{CLIR@10})=+0.96$, robust to dropping
the two collapsed encoders), so the lever is representation alignment at
indexing time, not a monolingual re-ranker at query time. On a cost-vs-capability
frontier, \texttt{embeddinggemma} is the Pareto-optimal capability corner (not the
cheapest reader), and machine-translating the \emph{question} is safe (no
significant penalty), so human-translation budget belongs on the corpus.
\end{abstract}

% ===========================================================================
\section{Introduction}
% ===========================================================================

A chemistry-patent search team that goes multilingual faces a deceptively
simple engineering decision: pick \emph{one} dense multilingual embedding model
to index a cross-jurisdictional patent corpus, and ship it. The dashboard that
guides that choice almost always shows a single number---mean Recall@10 on a
held-out query set---and the team picks the model on top. We argue that this
number is the wrong number, and that it systematically hides the failure the
team will actually deploy.

The reason is structural. In a multilingual patent corpus, the same invention
appears in several languages, and a query in one language must routinely reach
a document in another (a German chemist searching an English or Chinese family
member). Average recall, however, is dominated by the \emph{easy}
same-language case: when a same-language copy of the answer exists, every model
finds it. The hard, deployment-relevant case---retrieving a \emph{foreign}
language version---is averaged away. Our cross-lingual benchmark makes this
concrete by construction: it contains 34 Spanish queries with \emph{zero}
Spanish gold documents, a built-in ``no-home'' stress test that average recall
cannot see.

Two distinct failures hide behind that average, and both are deployment bugs.
First, \textbf{inconsistency}: ask the same compound in five languages and the
five ranked lists barely overlap---the best cross-lingual rank-biased-overlap
(RBO) \emph{any} of our nine models achieves is only $0.39$ on the alias-graph
benchmark and $0.19$ on the cross-lingual benchmark, meaning two language
variants of a query share fewer than four of their top-10 documents. Second,
\textbf{confusion}: a chemically-confusable \emph{wrong} compound out-ranks
every correct patent on $14$--$78\%$ of queries depending on the model. A search
team that reports only mean recall sees neither. And both are \emph{costly}: at
the median, finding a foreign twin takes scanning $\sim$3.5$\times$ more
documents than a same-language copy (\S\ref{sec:results}), the deficit lives
at the embedding level---so the fix is alignment, not re-ranking---and the model
choice is a cost-vs-capability frontier, not a single cheapest model.

We build the instruments to measure both failures, quantify what they cost,
identify the model that survives them, and trace the cause to the embedding
level. Our contributions are:

\begin{itemize}
  \item \textbf{C1.} The first \emph{content-controlled},
  chemistry-ontology-grounded cross-lingual patent-retrieval benchmarks, built
  only from human-translated patent text: an \emph{alias-graph} benchmark whose
  gold is defined by a chemical ontology and whose hard negatives are confusable
  neighbours, and a \emph{cross-lingual} benchmark over a 23{,}487-document
  shared corpus whose gold is genuinely parallel human-translated patents---not
  \texttt{publication\_number} equivalence and not machine-translated documents
  (\S\ref{sec:benchmarks}).
  \item \textbf{C2.} A cross-lingual robustness-metric family reported
  co-equally with recall---CLIR@k and the home advantage gap, the directional
  CLIR matrix (CLIRMatrix-style), mate-retrieval, cross-lingual RBO,
  language-collapse, and embedding separability AUC---anchored by two new
  deployment-cost instruments: \emph{XRC}, the cross-lingual reading-cost
  multiplier, presented on a capability-conditioned $(\mathrm{XRC50},
  \mathrm{CLIR@10})$ Pareto frontier, and \emph{RRC}, the re-ranker
  recoverability ceiling, read as a per-model budget curve with a knee $K^{\!*}$
  and a structural floor $L_\infty$, which an \emph{alignment-recoverability
  index} (ARI) decomposes into a re-ranker-recoverable, a deeper-pool, and an
  alignment-only share that sum to one. A reproducible degeneracy gate
  ($\mathrm{CLIR@10}<0.10$) defines which models are excluded; we use the
  CLIR-MRS / MRS composite only to order table rows, not as a contribution
  (\S\ref{sec:metrics}).
  \item \textbf{C3.} A mechanism finding, confirmed on a content-controlled
  corpus and made falsifiable: availability sets the stage but a residual
  encoder bias remains, the modal confusion is a same-language sibling, and the
  collapse is an embedding-level separability deficit---so a monolingual
  re-ranker cannot recover under-scored foreign twins, a claim we bound
  per-model with RRC (\S\ref{sec:analysis}).
  \item \textbf{C4.} A concrete deployment recommendation with an operating
  rule, grounded in a capability-conditioned cost frontier and a re-ranker-budget
  knee rather than mean recall or a hand-weighted composite: deploy
  \texttt{embeddinggemma} (the Pareto-optimal capability corner, with
  \texttt{bge-m3} the cheaper-to-read alternative), report
  XRC/RRC/CLIR@10/language-parity next to recall, budget the re-ranker by the
  knee, do not reflexively ensemble, and spend human-translation budget on the
  corpus rather than the queries; a per-route frontier shows the single best
  model is route-dependent, so a per-language router is genuine upside headroom
  rather than a delivered win (\S\ref{sec:deployment}).
  \item \textbf{C5.} A reproducible QAC generation and audit pipeline the
  benchmarks rest on, with human validation summarized in the system
  description (\S\ref{sec:benchmarks}, Appendix).
\end{itemize}

In one line: report robustness, not mean recall; and fix cross-lingual
retrieval by aligning representations at index time, not by re-ranking at query
time. Figure~\ref{fig:teaser} previews the gap that motivates the paper---every
model loses ground when no same-language document exists.

\begin{figure}[t]
  \centering
  \includegraphics[width=\linewidth]{cp_fig01_clir_leaderboard.png}
  \caption{Cross-lingual vs.\ same-language retrieval, nine models. Overall
  Recall@10 (grey) sits between the easy same-language case (MoLIR@10, green)
  and the hard cross-lingual case (CLIR@10, red). The best cross-lingual
  Recall@10 is only $0.50$ (\texttt{embeddinggemma}); average recall hides the
  collapse. Source: chem-patents key findings, Fig.~1.}
  \label{fig:teaser}
\end{figure}

% ===========================================================================
\section{Related Work}
% ===========================================================================

\paragraph{Multilingual and cross-lingual retrieval benchmarks.}
MIRACL \citep{miracl2023}, MMTEB \citep{mmteb2025}, and the TREC NeuCLIR track
\citep{neuclir2023} established large multilingual and cross-lingual retrieval
evaluations, and MTEB \citep{muennighoff2023mteb} is the de facto leaderboard
practitioners consult. CLIRMatrix \citep{clirmatrix2020} is the canonical
query-language~$\times$~document-language matrix and the precedent for the
directional analysis we apply; we use its style of decomposition on parallel
patents rather than Wikipedia. These efforts report a \emph{single averaged}
score per model; our contribution is orthogonal---a \emph{retrieval-side,
ranking-level} robustness suite reported co-equally with recall, purpose-built
to expose same-language bias that an average conceals. Because such rankings are
themselves sensitive to the aggregation scheme \citep{rankingrobustness2026} and
the community is moving toward consistency-aware aggregation (e.g.\ Borda in
MMTEB \citep{mmteb2025}), we test our recommendation across four aggregation
schemes and lean on per-axis dominance rather than on a single composite
(\S\ref{sec:results}). Recent work on benchmark contamination and maintenance
\citep{mteb_contamination2025} further motivates a fresh, content-controlled
evaluation asset.

\paragraph{Cross-lingual RAG and language preference.}
BordIRlines \citep{bordirlines2024}, XRAG \citep{xrag2025}, and Linguistic
Nepotism \citep{nepotism2025} document that retrievers and RAG systems prefer
documents in particular languages. Crucially, those studies confound
\emph{content} with \emph{language}: their multilingual documents differ in
content (Wikipedia articles that are not translations of one another) or are
machine-translated English imports (translationese as a confound). Our parallel
\emph{human-translated} patents are the content-controlled asset that isolates
the language effect from the content effect---the defense for C1.

\paragraph{Cross-lingual ranking: alignment, not translation.}
Two recent studies on general benchmarks reach a thesis close to ours:
\emph{What Drives Cross-Lingual Ranking?} \citep{whatdrivesclir2025} attributes
the cross-lingual gap to weak semantic alignment rather than translation and
argues systems should prioritise alignment over translation pipelines, and
\emph{The Cross-Lingual Cost} \citep{crosslingualcost2025} documents the
same-language head start and directional asymmetry. We do not claim to discover
the alignment-not-translation finding; we \emph{confirm} it on a
content-controlled parallel patent corpus that removes the translationese and
content confounds those studies could not, and add a chemistry-specific
confusability trap and a falsifiable per-model re-ranker bound (RRC).

\paragraph{Calibration, conformal, and fairness in IR.}
A growing line scores the \emph{calibration} of retrieval and RAG pipelines:
TRAQ \citep{traq2023}, CONFLARE \citep{conflare2024}, and Conformal-RAG
\citep{conformalrag2025} apply conformal prediction, while
fairness-via-optimal-transport post-processing
\citep{fairnessot2023,buyl2022optimal} re-balances scores. Our suite is
complementary but distinct: it scores \emph{ranking robustness} (cross-lingual
RBO, confusion, separability), not calibration. A reviewer familiar with these
tools should read our metrics as measuring \emph{whether the right foreign
document is ranked highly and consistently}, a question conformal coverage does
not answer. A finite-sample conformal version of our XRC cost metric---a
distribution-free \emph{guarantee} on cross-lingual reading depth, for which
Conformal-RAG \citep{conformalrag2025} is the machinery---is a natural horizon
we flag but do not claim here; the empirical XRC we report stands on its own.

\paragraph{Patent IR and patent-family non-equivalence.}
Patent retrieval has a long history (e.g.\ the PRIME line of NTCIR patent
experiments \citep{prime2002}), and CLEF-IP \citep{clefip2013} is the canonical
\emph{multilingual} (en/de/fr) patent-retrieval benchmark. Multilingual patent
retrieval thus already exists---but CLEF-IP's gold is prior-art relevance
judgments, not parallel human translation, and it is not chemistry/ontology
grounded; we narrow our ``first'' against it on exactly those axes. A key reason
for our benchmark design is that members of a patent family ``may differ in
their claims,'' so \texttt{publication\_number}-equivalence is an unsafe gold
signal; DAPFAM \citep{dapfam2025} deliberately uses family-level gold to handle
cross-jurisdiction redundancy---the design we reject. We instead define gold by
ontology membership (alias-graph) and by genuinely parallel human translation
(cross-lingual benchmark), never by family-number equivalence.

\paragraph{Chemistry IR and entity representations.}
SapBERT \citep{sapbert2021} aligns biomedical entity representations and PaECTER
\citep{paecter2024} learns patent-level representations, but both are largely
monolingual. A concurrent study \citep{patentembeddings2026} benchmarks
EmbeddingGemma, BGE-M3, and Qwen3 on \emph{English-only} patent retrieval; we
rank the same family on \emph{cross-lingual} robustness, where the ordering
differs. We evaluate cross-lingual chemistry retrieval with ontology-graph
grounding (ChEBI, \citealp{chebi2016}) and confusable-neighbour hard negatives.

\paragraph{When to translate.}
Whether to translate the query or the documents is a 25-year-old CLIR question
\citep{oard1998,salehpecina2020}, with query translation often sufficing in a
specialised domain \citep{salehpecina2020}; pre-neural CLIR also studied
translation-quality prediction and routing \citep{kishida2008}. Our budget rule
(\S\ref{sec:deployment}) re-derives this for embedding retrieval over patents
and quantifies the query-side MT cost as an insignificant null. For
cross-lingual RBO we treat the five language versions of a query as query
\emph{variants} and apply retrieval-consistency analysis in the lineage of
\citet{bailey2017consistency}, cross-lingually.

Having positioned our four contributions against CLEF-IP/DAPFAM, the
alignment-not-translation line, and the QT-vs-DT literature, we now build the
benchmarks that deliver C1---and the two design choices (no MT source
documents, no \texttt{publication\_number} gold) that the patent-IR and
cross-lingual-RAG lines above show are necessary.

% ===========================================================================
\section{Benchmarks}
\label{sec:benchmarks}
% ===========================================================================

We release two benchmarks, both grounded in real chemistry-patent text and both
deliberately avoiding the two shortcuts that make a ``multilingual'' benchmark
misleading: machine-translated source documents (translationese) and
family-number gold (non-equivalence).

\paragraph{Shared corpus.}
Both benchmarks retrieve against one shared haystack, the
\texttt{multilingual\_GP} corpus of \textbf{23{,}487} documents drawn from
Google Patents and EPO bulk data in English, German, Spanish, French, and
Chinese, so that results are directly comparable across benchmarks.\footnote{%
Corpus row count, chem-patents key findings header.} Detailed
corpus-construction statistics (source counts after de-duplication; IPC/A61
distribution) are deferred to the system description; all load-bearing sizes
below come from the two benchmark datasets.
% TODO trace: corpus dedup count 14,401 and GP/EPO/JRC coverage matrix live in
% the source system slides, not yet under reports/; restrict load-bearing
% numbers to the two key_findings files until dumped.

\paragraph{Alias-graph benchmark.}
Here a query is a multilingual ChEBI alias set for a compound (e.g.\ tin:
\emph{tin}/Sn, \emph{Zinn}, \emph{esta\~no}, \emph{\'etain}), the gold set is the
patents about that concept, and the hard negatives are documents about
chemically-confusable neighbours (siblings, parents, tautomers) taken from the
compound's alias graph. The benchmark has \textbf{132} questions over
\textbf{24} ChEBI compounds in five languages, and reports \emph{two co-equal
relevance lenses}: a \emph{concept} lens (any patent about the compound;
$\sim$109 gold per query) and a \emph{per-publication} lens (the query's own
source patent and its translations; $\sim$2.4 gold per query).\footnote{Alias
benchmark size and lens definitions, alias-graph key findings header.} The
per-publication lens is the one we read as a quality score; the concept lens is
capped by its $\sim$109 positives and must not be read as recall quality.

\paragraph{Cross-lingual (CLIR) benchmark.}
Here a query is a generated question about a source patent, and the gold set is
\emph{every language version} of that source patent. The benchmark has
\textbf{137} questions in English, German, Spanish, French, and Chinese, split
into two non-symmetric populations: \textbf{57 original} questions, asked in the
source patent's own language (so each has exactly one same-language gold), and
\textbf{80 synthetic} questions, machine-translated into another language (so
they have \emph{no same-language gold at all}---pure cross-lingual
retrieval).\footnote{Population split and Spanish property, chem-patents key
findings header and ``defining property.''} \textbf{Spanish is a pure
query-side language}: 34 Spanish queries map to zero Spanish gold documents,
giving us a controlled ``no-home'' setting in which average recall has nowhere
easy to hide.\footnote{Benchmark defining property, chem-patents key findings.}

\paragraph{Honesty by design.}
Two choices keep the benchmarks honest. (i) \emph{Source text is
human-translated only.} We never use machine-translated documents in the
corpus, so a cross-lingual miss is a retrieval failure, not a translationese
artifact. (ii) \emph{Gold is defined by ontology or by human translation, never
by \texttt{publication\_number}}, side-stepping the patent-family
non-equivalence problem (\S2). Machine translation is used only to translate
\emph{questions} for the synthetic population, and \S\ref{sec:results} shows
that this introduces no significant cross-lingual penalty.

\paragraph{Generation and validation pipeline (C5).}
The queries are produced by a reproducible per-(document, language, mode)
generation pipeline that emits three scored candidates best-first, graded on
faithfulness and answerability, across technical and semantic question modes. A
human annotator reviewed a sample of the generated items; the human-rated
quality and the LLM auto-grader's strictness relative to humans are summarized
qualitatively in the system description. We keep C5 as supporting evidence: the
benchmarks' validity rests on the design choices above, not on these validation
numbers.
% TODO trace: human-eval numbers (e.g. mean 8.33/10, 97/100 reviewed, 0
% rejected, auto-grader +4.3pp stricter) live in the source system slides, not
% yet under reports/; either restrict to descriptive text or dump to reports/.

% ===========================================================================
\section{Metrics}
\label{sec:metrics}
% ===========================================================================

Recall measures \emph{capability}; it does not measure \emph{robustness to
language}. We define a family of metrics that isolate the cross-lingual axes a
single recall number averages away. All are computed per model on both
benchmarks.

\paragraph{CLIR@k, MoLIR@k, and home advantage.}
We split Recall@k by whether the retrieved gold is in the query's own language.
\emph{MoLIR@k} (mono-lingual) restricts to same-language gold; \emph{CLIR@k}
restricts to cross-language gold. The \emph{home advantage}
$\mathrm{MoLIR}{-}\mathrm{CLIR}$ is the size of the same-language head start.
This pair is the direct instrument for ``average recall hides collapse.''

\paragraph{Directional CLIR and asymmetry.}
Following the CLIRMatrix-style decomposition \citep{clirmatrix2020}, we compute a
query-language $\rightarrow$ document-language recall matrix and its asymmetry,
because retrieval is not a symmetric similarity: the recall of reaching language
$b$ from a query in language $a$ need not equal the reverse. We read it as a
\emph{hub-and-spoke graph} (which languages are easy targets, which directed
edges are hardest, which pairs are most asymmetric).

\paragraph{Mate-retrieval.}
For a query whose answer also exists as a foreign \emph{twin} (a parallel
translation), we measure whether that twin is ever retrieved and at what depth:
mate-hit@k, mate-MRR, and the rank of the first foreign twin. This isolates
``can the foreign version be found at all'' from ``is it ranked first.''

\paragraph{Cross-lingual RBO.}
For a compound asked in five languages, we measure rank-biased overlap
\citep{rbo2010} between the five top-$k$ lists; RBO$=1$ means identical
rankings, $0$ means disjoint. This is our \emph{consistency} axis: does the
deployed index return the same patents regardless of query language. We treat
the five language versions as query \emph{variants} and apply RBO
cross-lingually, in the retrieval-consistency-under-query-variation lineage of
\citet{bailey2017consistency}.

\paragraph{Language collapse / over-representation.}
We measure how much a query language over-fetches documents in its own
language relative to the corpus base rate, and how often same-language
\emph{noise} out-ranks foreign \emph{gold}. This is the bias mechanism behind
the home advantage.

\paragraph{Separability AUC.}
ROC-AUC over gold-vs-non-gold similarity scores is a standard separability
diagnostic; our use is the same-vs-cross \emph{decomposition} and the re-ranker
corollary. We measure $\mathrm{AUC}(\text{score(gold)} > \text{score(non-gold)})$
separately for same-language and cross-language gold. A low cross-language AUC
means foreign gold is \emph{under-scored}, not merely mis-ordered---the
distinction that decides whether a re-ranker can help.

\paragraph{XRC: cross-lingual reading-cost multiplier.}
Separability tells us foreign gold is under-scored; XRC turns that into a
deployment-legible cost. Let $D_C^{\text{same}}$ and $D_C^{\text{cross}}$ be the
retrieval depths at which a coverage fraction $C$ of same-language and
cross-language gold (resp.\ foreign twins) is reached. Then
\begin{equation}
\label{eq:xrc}
\mathrm{XRC}_C(m) = D_C^{\text{cross}}(m) \,/\, D_C^{\text{same}}(m),
\end{equation}
the factor by which a practitioner must read deeper to reach a foreign twin than
a same-language copy. Unlike a unitless score, XRC is in \emph{documents read}
and is monotone-invariant (a ratio of retrieval depths is invariant to any
monotone re-scaling of similarities, unlike an AUC or a weighted composite).
Here $D^{\text{same}}_C$ is taken over the $57$ same-language-gold queries and
$D^{\text{cross}}_C$ over the $137$ cross-gold queries, so XRC is a
\emph{population-level}, not a paired within-query, ratio. We report
\textbf{XRC50} (the median, $C{=}0.5$) as the headline because it is finite at
our sample size; the deeper coverages $D_{90}/D_{95}$ are right-censored (a
fraction of first-foreign ranks is infinite) and are reported only as lower
bounds, never as point estimates.

\paragraph{RRC: re-ranker recoverability ceiling.}
A re-ranker can only reorder a candidate pool it is given. We define
\begin{equation}
\label{eq:rrc}
\mathrm{RRC@}K(m) = \Pr\!\left[\text{first foreign twin rank} \le K\right],
\end{equation}
the fraction of cross-lingual queries whose foreign twin appears within the
top-$K$ candidates; $1-\mathrm{RRC@}K$ is then \emph{provably unrecoverable} by
any top-$K$ re-ranker. RRC@$K$ is the cumulative first-foreign-twin hit rate
(mate-hit@$K$ restricted to cross-lingual queries); our contribution is not the
quantity but reading it as a per-model re-ranker \emph{budget curve}---the knee
$K^{\!*}$ past which a deeper candidate pool buys almost nothing, and the floor
$L_\infty = 1-\mathrm{RRC@}K_{\max}$ that no re-ranker over the top-$1000$ can
move. The knee and the diminishing-returns-with-depth shape are the established
multi-stage cascade result that recall is capped by first-stage candidate depth
\citep{nogueira2019multistage,gao2021rethink}; our contribution is its
cross-lingual quantification and the structural floor $L_\infty$. RRC thus
converts ``align, don't re-rank'' from a slogan into a per-model upper bound with
an explicit structural floor.

\paragraph{ARI: the alignment-recoverability decomposition.}
The RRC curve admits a natural exhaustive reading of the same shortfall. At any
depth $K$, a model's cross-lingual shortfall splits into a cheaply recoverable
part ($\mathrm{RRC@}K$, what a top-$K$ re-ranker can reach), a deep-pool part
($\mathrm{RRC@}1000-\mathrm{RRC@}K$, recoverable only by a deeper candidate
pool), and the alignment-only floor $L_\infty=1-\mathrm{RRC@}1000$; the three sum
to one. We report the scalar
\begin{equation}
\label{eq:ari}
\mathrm{ARI@}K(m) = \frac{L_\infty(m)}{1-\mathrm{RRC@}K(m)},
\end{equation}
the \emph{alignment-recoverability index}: the fraction of the \emph{remaining}
cross-lingual gap at depth $K$ (after a top-$K$ re-rank) that no re-ranker can
recover and only representation alignment can move. Normalizing a re-ranking
remainder by a recoverable gap has a precedent in long-tailed reranking
\citep{residualrerank2026}; we invert it (the un-rerankable share rather than the
reranker's recoverable gain) and tie it to representation alignment
cross-lingually, with an alignment-only floor that $\rho_k$ has no analogue for.
ARI is the quantitative form of ``align, don't re-rank''---a per-model split that
closes to one (Figure~\ref{fig:ari}). Source: \texttt{extra\_ari\_decomposition}.

\paragraph{The degeneracy gate.}
Several axes (XRC, RRC, separability) are undefined or meaningless for a model
that retrieves almost nothing, and we exclude such models from every
``non-degenerate'' summary. We make the exclusion reproducible with a single
criterion: a model is \emph{degenerate} iff $\mathrm{CLIR@10} < 0.10$. This gate
flags \emph{exactly} \texttt{gte-base} and \texttt{e5-large-instruct}
(Figure~\ref{fig:cp_deg}), matching the two models the rest of the paper
excludes; every later ``non-degenerate'' qualifier is anchored to it.\footnote{A
stricter two-criterion gate ($\mathrm{CLIR@10}<0.10$ \emph{and}
$\mathrm{RRC@}1000<0.10$) flags only \texttt{gte-base}, because
\texttt{e5-large-instruct} surfaces enough twins deep in the list
($\mathrm{RRC@}1000=0.63$) to clear the second criterion; the single CLIR@10
criterion is the clean rule that recovers both exclusions. Source:
\texttt{extra\_two\_tax\_degeneracy}.}

\begin{figure}[t]
  \centering
  \includegraphics[width=\linewidth]{cp_fig20_degeneracy_gap.png}
  \caption{The degeneracy gate $\mathrm{CLIR@10}<0.10$ (dashed line) separates
  \texttt{gte-base} (CLIR@10 $=0.000$, bar not visible) and
  \texttt{e5-large-instruct} ($0.077$) from the seven non-degenerate models
  (lowest \texttt{SapBERT} $0.179$). Source: chem-patents
  \texttt{extra\_two\_tax\_degeneracy}.}
  \label{fig:cp_deg}
\end{figure}

\paragraph{CLIR-MRS / MRS (table-ordering convenience only).}
For completeness we report a composite. On the cross-lingual benchmark,
\begin{equation}
\label{eq:clirmrs}
\mathrm{CLIR\text{-}MRS} = \mathrm{cap} \times \bigl(0.5 + 0.5\cdot\mathrm{rob}\bigr),
\end{equation}
where capability $\mathrm{cap}$ aggregates accuracy, CLIR, and separability and
robustness $\mathrm{rob}$ aggregates consistency, MT-robustness, and
language-parity ($\pm50\%$ modulation).\footnote{CLIR-MRS definition,
chem-patents key findings ``Verdict.''} The alias-graph MRS is the
min--max-normalised mean of accuracy, consistency, confusion-robustness,
language-parity, and separability, with a 95\% bootstrap CI.\footnote{MRS
definition, alias-graph key findings ``Headline numbers'' footnote.} We use
these composites \emph{only to order table rows}; the per-axis numbers carry the
argument, and \S\ref{ssec:boards} shows the model ranking is aggregation-sensitive,
so no composite weighting is load-bearing.

% ===========================================================================
\section{Experimental Setup}
\label{sec:setup}
% ===========================================================================

\paragraph{Models.}
We evaluate \textbf{nine} multilingual embedding models as drop-in dense
retrievers: \texttt{embeddinggemma}, \texttt{bge-m3} \citep{bgem3_2024},
\texttt{granite-278m}, \texttt{nomic-v2-moe}, \texttt{qwen3-0.6B},
\texttt{LaBSE} \citep{labse2022}, \texttt{SapBERT} \citep{sapbert2021},
\texttt{e5-large-instruct}, and \texttt{gte-base}.

\paragraph{Protocol.}
Both benchmarks retrieve against the single shared \texttt{multilingual\_GP}
haystack (23{,}487 documents); queries and qrels come from the respective
benchmark datasets. We report Recall@10 and the full metric family of
\S\ref{sec:metrics}. For the alias-graph benchmark we report both relevance
lenses; for the cross-lingual benchmark we report the original/synthetic split.

\paragraph{Reproducibility.}
Every figure and number is regenerated by
\texttt{reports/runs/chem\_patents/experimental\_codes/run\_all.py} and the
corresponding \texttt{reports/runs/alias\_graph/experimental\_codes/} scripts,
which also emit the executive summaries this paper cites. The XRC/RRC,
directional-hub, aggregation-ribbon, availability-residual, joint-failure, and
confusion-severity analyses---now including the cost frontier, RRC budget
frontier, two-tax/degeneracy-gate, ARI decomposition, and per-route frontier
analyses (\texttt{extra\_cost\_frontier.py},
\texttt{extra\_rrc\_budget\_frontier.py},
\texttt{extra\_two\_tax\_degeneracy.py}, \texttt{extra\_ari\_decomposition.py},
\texttt{extra\_per\_route\_frontier.py})---are regenerated by the
\texttt{experimental\_plots/extra\_*.py} scripts under the two run families.

% ===========================================================================
\section{Results}
\label{sec:results}
% ===========================================================================

We report the cross-lingual benchmark findings first (\S\ref{ssec:cp}), then
the alias-graph findings (\S\ref{ssec:ag}), then the two leaderboards
(\S\ref{ssec:boards}).

\subsection{The cross-lingual collapse}
\label{ssec:cp}

\paragraph{Collapse is real and large.}
Cross-lingual recall trails same-language recall in \emph{every} model
(Figure~\ref{fig:teaser}). The best CLIR@10 is $0.50$ (\texttt{embeddinggemma}),
and the home advantage $\mathrm{MoLIR}{-}\mathrm{CLIR}$ reaches $+0.55$ for the
most biased model (\texttt{e5-large-instruct})---much of which gold availability
shapes (\S\ref{sec:analysis}), though a residual encoder bias remains: a
same-language copy of the answer is far easier to retrieve than its foreign twin
(Figure~\ref{fig:cp_home}).\footnote{MoLIR@10 is defined only on the 57 original
queries (the only ones with a same-language gold); the home advantage is
measured \emph{paired} within those queries, so the $+0.55$ is a clean
within-query contrast even though the green MoLIR bar in Figure~\ref{fig:teaser}
is a smaller population than the grey/red bars.}

\begin{figure}[t]
  \centering
  \includegraphics[width=\linewidth]{cp_fig02_home_advantage.png}
  \caption{Same-language home advantage ($\mathrm{MoLIR}{-}\mathrm{CLIR}$) per
  model on the cross-lingual benchmark; it reaches $+0.55$.
  Source: chem-patents key findings, Fig.~2.}
  \label{fig:cp_home}
\end{figure}

\paragraph{Retrieval is anisotropic: a hub-and-spoke graph.}
Read as a directed graph (Figure~\ref{fig:cp_dir}), the
query-language~$\rightarrow$~document-language matrix is not symmetric: the
hardest directed edge is en$\rightarrow$de (R@10 $=0.12$) and the most
asymmetric pair is de$\leftrightarrow$zh (gap $+0.23$). Contrary to a common
assumption, \emph{no single language is a clean ``easiest target''}: pooled over
the eight models with a defined cross-lingual recall (all but the degenerate
\texttt{gte-base}), French ($0.375$) and English ($0.367$) are
nearly tied as targets (within $0.01$ of each other), with Chinese ($0.350$) and
German ($0.309$) behind. What asymmetry exists partly tracks corpus composition (en
$46\%$ vs.\ zh $0.4\%$ of documents), not only encoder behaviour---there are
simply few German or Chinese documents to retrieve. A single ``recall'' number
erases all of this directionality.

\begin{figure}[t]
  \centering
  \includegraphics[width=\linewidth]{cp_fig03_directional_clir_matrix.png}
  \caption{Query$\rightarrow$document language recall matrix (best model);
  retrieval direction matters, but no language is a clean easiest target. The
  most asymmetric directed pair is de$\leftrightarrow$zh (gap $+0.23$; asymmetry
  panel not shown). Source: chem-patents key findings, Fig.~3; pooled hub scores
  from \texttt{extra\_directional\_hub}.}
  \label{fig:cp_dir}
\end{figure}

\paragraph{The cost of cross-linguality, and the cost frontier.}
How much does the collapse cost in practice? At the median, the best model
\texttt{embeddinggemma} reaches a same-language gold at depth $2$ but its foreign
twin only at depth $7$---an XRC50 of $\mathbf{3.5\times}$ (Eq.~\ref{eq:xrc}). The
cleanest model on this axis is \texttt{granite-278m} ($1.25\times$); the cost
balloons for \texttt{nomic-v2-moe} ($11.5\times$) and becomes catastrophic for
\texttt{e5-large-instruct} ($97.75\times$), while \texttt{gte-base} is degenerate
(\S\ref{sec:metrics}: $\mathrm{CLIR@10}<0.10$) and its XRC is undefined.

A low reading cost is worthless if the model retrieves nothing, so XRC is not a
scalar to minimize: we read it \emph{against} capability on the standard
cost-vs-capability Pareto frame \citep{syftr2025}, but with XRC---our new
reading-depth multiplier---as the cost axis in place of the usual latency or
parameter count, giving $(\mathrm{XRC50},\mathrm{CLIR@10})$ space
(Figure~\ref{fig:cost_frontier}). Three models are Pareto-optimal---no other
model is both cheaper-to-read \emph{and} higher-CLIR:
\texttt{embeddinggemma} (XRC50 $3.5\times$, CLIR@10 $0.50$), \texttt{bge-m3}
($2.0\times$, $0.44$), and \texttt{granite-278m} ($1.25\times$, $0.33$).
\texttt{embeddinggemma} is the frontier's \emph{unique maximum-CLIR corner}; it
is \emph{not} the cheapest deployable model. At a stated (untuned) deployment
threshold $\tau{=}0.40$ the admitted set is
$\{$\texttt{bge-m3}, \texttt{qwen3-0.6B}, \texttt{embeddinggemma}$\}$ and the
cheapest-to-read admitted model is \texttt{bge-m3} ($2.0\times$ vs.\ $3.5\times$),
not \texttt{embeddinggemma}. This snapshot survives only over a narrow band:
sweeping $\tau$, the admitted set is stable for $\tau \in [0.39, 0.43]$ and
\texttt{bge-m3} remains the cheapest admitted reader for $\tau \in [0.33, 0.44]$;
below $\tau{\approx}0.33$ the cheaper-to-read \texttt{granite-278m}
($1.25\times$) enters the admitted set and the cheapest-reader recommendation
\emph{flips} to \texttt{granite-278m}, while above $\tau{\approx}0.45$ only
\texttt{embeddinggemma} qualifies. Only \texttt{embeddinggemma}'s status as the
unique maximum-CLIR@10 corner is $\tau$-invariant. The picture also reads off a
cautionary point: the
cheapest reader \emph{can} be among the worst retrievers (the dominated
look-alikes like \texttt{SapBERT} read shallow precisely because they surface
little)---a directional read-off of the figure, not a statistical claim (the
XRC50$\sim$CLIR@10 rank correlation across the seven non-degenerate models is
$\rho{=}{+}0.29$, $n{=}7$, $p{=}0.53$, n.s.). The deployment consequence is a
\emph{frontier choice}, not a single cheapest model (\S\ref{sec:deployment}).

\begin{figure}[t]
  \centering
  \includegraphics[width=\linewidth]{cp_fig18_cost_frontier.png}
  \caption{Cost-vs-capability frontier in $(\mathrm{XRC50},\mathrm{CLIR@10})$
  space. The Pareto frontier (line) has three members
  \{\texttt{embeddinggemma}, \texttt{bge-m3}, \texttt{granite-278m}\};
  \texttt{embeddinggemma} is the unique maximum-CLIR corner. Dominated finite
  models carry an $\times$; \texttt{gte-base} is off-plane (XRC undefined).
  \texttt{embeddinggemma} is the capability corner, \emph{not} the cheapest
  reader---at $\tau{=}0.40$ that is \texttt{bge-m3} ($2.0\times$), stable for
  $\tau \in [0.33, 0.44]$ and flipping to \texttt{granite-278m} below
  $\tau{\approx}0.33$; only the \texttt{embeddinggemma} corner is
  $\tau$-invariant. Source: chem-patents \texttt{extra\_cost\_frontier}.}
  \label{fig:cost_frontier}
\end{figure}

\paragraph{Machine-translating the question is safe.}
Controlling for the source patent, the paired human$-$MT difference in
cross-lingual reach is $-0.044$ ($p=0.13$): statistically
\emph{insignificant} (Figure~\ref{fig:cp_mt}). We read this as a \emph{null}
result---MT of the query does not significantly hurt and we do not claim it
helps---which validates the project's ``MT-is-fine-for-the-question''
assumption.

\begin{figure}[t]
  \centering
  \includegraphics[width=\linewidth]{cp_fig05_mt_penalty.png}
  \caption{Paired human vs.\ machine-translated queries; the cross-lingual
  difference is $-0.044$ ($p=0.13$), not significant.
  Source: chem-patents key findings, Fig.~5.}
  \label{fig:cp_mt}
\end{figure}

\paragraph{Foreign twins are buried, and a re-ranker cannot reach them all.}
Pooled mate-hit@10 is $0.38$, and on $15\%$ of (query, model) pairs the foreign
twin never surfaces even in the top-1000 (Figure~\ref{fig:cp_mate}). The best
twin-finder, \texttt{embeddinggemma}, reaches its first foreign twin at a median
rank of $5$ (Figure~\ref{fig:cp_mate})---most models do far worse. This directly
bounds any downstream re-ranker (Eq.~\ref{eq:rrc}), and we read the bound as a
\emph{budget curve} (Figure~\ref{fig:rrc_budget}): for \texttt{embeddinggemma}
the RRC$(K)$ knee is $K^{\!*}{=}5$, so most of the recoverable payoff sits in a
shallow top-$K$ and a deeper pool past the knee buys almost nothing; the knee is
$\le 20$ for nearly every non-degenerate model. The ceilings are RRC@100 $=0.7445$
and RRC@1000 $=0.9416$, so a realistic top-100 re-ranker recovers at most $74\%$
of foreign twins (leaving $\sim$25\% on the table) and the floor
$L_\infty = 1-\mathrm{RRC@1000} = \mathbf{0.058}$ of foreign twins is
\textbf{unrecoverable by \emph{any} re-ranker over the top-1000}---a floor only
representation alignment can move (\S\ref{sec:analysis}). This floor ranges from
$0.058$ (\texttt{embeddinggemma}) to $0.372$ for the degenerate
\texttt{e5-large-instruct}, which the gate excludes; the regression checks behind
the curve all pass, so we state it firmly. Re-ranking reads a list retrieval
already truncated.

\paragraph{How much of the gap is alignment-bound: the ARI decomposition.}
The ARI decomposition (Figure~\ref{fig:ari}) reports this split per model. For
\texttt{embeddinggemma} the alignment-only floor is the smallest of any
non-degenerate model ($L_\infty=0.058$, next \texttt{qwen3-0.6B} at $0.073$).
On the post-re-rank residual itself \texttt{embeddinggemma} and \texttt{qwen3-0.6B}
are \emph{tied} for the lowest re-ranker-irreducible share---$\mathrm{ARI@}100 =
0.229$ and $0.233$ respectively (Eq.~\ref{eq:ari}), a gap of $0.004$ whose
$95\%$ resample interval $[-0.174, 0.176]$ straddles zero (Appendix~\ref{app:robust}).
The decomposition is what turns ``align, don't re-rank'' from a slogan into a
measured per-model remainder.

\begin{figure}[t]
  \centering
  \includegraphics[width=\linewidth]{cp_fig22_ari_decomposition.png}
  \caption{Alignment-recoverability decomposition: each model's cross-lingual
  shortfall split into the re-ranker-recoverable part ($\mathrm{RRC@}100$), the
  deeper-pool part ($\mathrm{RRC@}1000-\mathrm{RRC@}100$), and the alignment-only
  floor $L_\infty$. The three sum to $1.0$ for every model (the identity closes
  for all nine; the figure shows the seven non-degenerate).
  \texttt{embeddinggemma} has the smallest floor ($0.058$) and is tied with
  \texttt{qwen3-0.6B} for the lowest post-re-rank residual
  ($\mathrm{ARI@}100=0.229$ vs.\ $0.233$). Source: chem-patents
  \texttt{extra\_ari\_decomposition}.}
  \label{fig:ari}
\end{figure}

\begin{figure}[t]
  \centering
  \includegraphics[width=\linewidth]{cp_fig19_rrc_budget.png}
  \caption{Re-ranker budget curves: RRC$(K)$ vs.\ candidate depth $K$ (log scale)
  for the eight models with a defined RRC curve (all but \texttt{gte-base}, whose
  candidate pool is empty), with the knee $K^{\!*}$ (rings) past which
  a deeper pool buys almost nothing. \texttt{embeddinggemma} knees at
  $K^{\!*}{=}5$, reaches RRC@1000 $=0.942$, and leaves an alignment-only floor
  $L_\infty=0.058$ unrecoverable by any re-ranker. Source: chem-patents
  \texttt{extra\_rrc\_budget\_frontier}.}
  \label{fig:rrc_budget}
\end{figure}

\begin{figure}[t]
  \centering
  \includegraphics[width=\linewidth]{cp_fig06_07_mate.png}
  \caption{Foreign twins are buried (cross-lingual benchmark). \emph{Left:}
  mate-retrieval (mate-hit@10/100, mate-MRR); pooled mate-hit@10 is $0.38$.
  \emph{Right:} depth to the first foreign twin; $15\%$ of (query, model) pairs
  never reach one in the top-1000. Source: chem-patents key findings,
  Figs.~6--7.}
  \label{fig:cp_mate}
\end{figure}

\subsection{Inconsistency and confusion (alias-graph)}
\label{ssec:ag}

If \S\ref{ssec:cp} measured what cross-linguality costs to \emph{read} (the XRC
reading-cost tax), the alias-graph benchmark measures what it costs in
\emph{precision}---the second line-item of the same bill: a look-alike compound
that out-ranks the gold (a confusability tax). The two taxes are only weakly---and
if anything inversely---rank-correlated across the seven non-degenerate models
(Spearman $\rho{=}{-}0.59$, $n{=}7$, $p{=}0.16$, n.s.; Figure~\ref{fig:two_tax}),
so neither benchmark is a clean proxy for the other. We read this descriptively,
as the motivation for reporting both benchmarks rather than one: a model can read
cheaply yet confuse heavily, or vice versa (\texttt{embeddinggemma} is low on both;
\texttt{granite-278m} reads cheapest at $1.25\times$ but confuses at $0.182$). It
is not an independence result.

\begin{figure}[t]
  \centering
  \includegraphics[width=\linewidth]{cp_fig21_two_tax.png}
  \caption{The two line-items of the cross-lingual tax per model: the reading-cost
  tax (XRC50, cross-lingual benchmark) against the confusability tax
  (confusion rate, alias-graph benchmark). Across the seven non-degenerate
  models the two are only weakly and inversely rank-correlated
  ($\rho{=}{-}0.59$, $n{=}7$, $p{=}0.16$, n.s.), so neither is a proxy for the
  other. Source: chem-patents \texttt{extra\_two\_tax\_degeneracy}.}
  \label{fig:two_tax}
\end{figure}

\paragraph{Same question, different rankings.}
Asking the same compound in five languages does \emph{not} return the same
patents. The best cross-lingual RBO \emph{any} of the nine models reaches is only
$0.39$ (Figure~\ref{fig:ag_rbo})---a ceiling no model beats; two language
variants of a query share fewer than four of their top-10 documents, and Chinese
is the consistent odd-one-out. The
embedding space is far from language-agnostic.

\begin{figure}[t]
  \centering
  \includegraphics[width=\linewidth]{ag_fig1_cross_lingual_rbo.png}
  \caption{Same compound, five languages: cross-lingual RBO. The ceiling is
  $0.39$ (\texttt{embeddinggemma}). Source: alias-graph key findings, Fig.~1.}
  \label{fig:ag_rbo}
\end{figure}

\paragraph{Confusable look-alikes out-rank the gold.}
A chemically-confusable wrong compound out-ranks every gold patent on
$14$--$78\%$ of queries (per-publication lens) depending on the model, worst for
German and Chinese (Figure~\ref{fig:ag_conf}). The damage is concentrated in
near-twins: pooled, a \emph{sibling} compound out-ranks the gold $18.1\%$ of the
time versus only $6.2\%$ for a more distant \emph{parent} class (a $2.9\times$
ratio; for \texttt{embeddinggemma}, $6.1\%$ vs.\ $1.5\%$).\footnote{This is a
two-level sibling/parent split; the on-disk relation field is binary, so a
graded ChEBI hop-distance decay law is left to future work. Source:
\texttt{extra\_confusion\_severity}.} The repeat offenders are a small set of
universal attractors---\emph{polypeptide}, \emph{methyl}, \emph{ethene},
\emph{hydroxide}, \emph{dioxygen} (Figure~\ref{fig:ag_attr}).

\begin{figure}[t]
  \centering
  \includegraphics[width=\linewidth]{ag_fig2_confusion_both_lenses.png}
  \caption{Confusion rate (a look-alike beats all gold), model $\times$ language,
  both lenses; $14$--$78\%$ on the publication lens.
  Source: alias-graph key findings, Fig.~2.}
  \label{fig:ag_conf}
\end{figure}

\begin{figure}[t]
  \centering
  \includegraphics[width=\linewidth]{ag_fig5_universal_attractors.png}
  \caption{Universal attractors: look-alikes that most often steal the top rank.
  Source: alias-graph key findings, Fig.~5.}
  \label{fig:ag_attr}
\end{figure}

\subsection{Leaderboards}
\label{ssec:boards}

Table~\ref{tab:cp_board} ranks the nine models by CLIR-MRS on the cross-lingual
benchmark and Table~\ref{tab:ag_board} by MRS on the alias-graph benchmark.
\texttt{embeddinggemma} tops both. Note the order changes from
Figure~\ref{fig:teaser}: ranking by CLIR-MRS rather than raw recall reshuffles
the middle of the field (\texttt{granite-278m} and \texttt{qwen3-0.6B} swap),
which is the point of the paper---robustness reorders the recall ranking
(\texttt{embeddinggemma} leads consistency and separability, not raw recall
alone; the per-axis detail is in Tables~\ref{tab:cp_board}--\ref{tab:ag_board}
and Figure~\ref{fig:cp_ribbon}).

\paragraph{The ranking is aggregation-sensitive.}
We caution against over-reading the composite. Across four aggregation schemes
(CLIR-MRS, winner-take-all, Borda, and equal-weight), \texttt{embeddinggemma}'s
rank \emph{ranges} from $1$ to $4$: it is rank-1 only when capability is
up-weighted, because it leads the three capability axes (accuracy, CLIR,
separability) but trails on the robustness axes
(Figure~\ref{fig:cp_ribbon}).\footnote{The winner-take-all (axes-won) column is
contaminated by a degenerate model: \texttt{gte-base} ``wins'' MT-robustness and
language-parity only because it retrieves almost nothing
($\mathrm{CLIR@10}=0.000$, below the degeneracy gate of \S\ref{sec:metrics}), so
its scores are trivially uniform.}
Rankings are therefore aggregation-sensitive; our recommendation
(\S\ref{sec:deployment}) rests on \emph{per-axis dominance}, not on any single
composite weighting.

\begin{figure}[t]
  \centering
  \includegraphics[width=\linewidth]{cp_fig17_aggregation_ribbon.png}
  \caption{Model rank under four aggregation schemes.
  \texttt{embeddinggemma}'s rank ranges $[1,4]$; the ordering is
  aggregation-sensitive. Source: chem-patents
  \texttt{extra\_aggregation\_invariance}.}
  \label{fig:cp_ribbon}
\end{figure}

\begin{table}[t]
  \centering
  \small
  \begin{tabular}{lrrrrr}
    \toprule
    model & R@10 & CLIR@10 & home & sep & CLIR-MRS \\
    \midrule
    \texttt{embeddinggemma} & 0.544 & 0.502 & $+0.22$ & 0.90 & \textbf{0.71} \\
    \texttt{bge-m3}         & 0.476 & 0.437 & $+0.23$ & 0.87 & 0.70 \\
    \texttt{granite-278m}   & 0.359 & 0.329 & $+0.17$ & 0.86 & 0.67 \\
    \texttt{nomic-v2-moe}   & 0.443 & 0.381 & $+0.34$ & 0.86 & 0.67 \\
    \texttt{qwen3-0.6B}     & 0.468 & 0.433 & $+0.19$ & 0.86 & 0.58 \\
    \texttt{LaBSE}          & 0.277 & 0.238 & $+0.21$ & 0.81 & 0.40 \\
    \texttt{SapBERT}        & 0.212 & 0.179 & $+0.18$ & 0.79 & 0.40 \\
    \texttt{e5-large-instruct} & 0.178 & 0.077 & $+0.55$ & 0.71 & 0.26 \\
    \texttt{gte-base}       & 0.004 & 0.000 & $+0.02$ & 0.49 & 0.00 \\
    \bottomrule
  \end{tabular}
  \caption{Cross-lingual benchmark leaderboard. ``home'' is the home advantage
  $\mathrm{MoLIR}{-}\mathrm{CLIR}$; ``sep'' is separability AUC.
  \texttt{embeddinggemma} wins (CLIR-MRS $0.71$, CI $[0.67,0.77]$).
  Source: chem-patents key findings, Verdict table / \texttt{headline\_numbers.csv}.}
  \label{tab:cp_board}
\end{table}

\begin{table}[t]
  \centering
  \small
  \begin{tabular}{lrrrr}
    \toprule
    model & pubR@10 & RBO & conf & MRS \\
    \midrule
    \texttt{embeddinggemma} & 0.67  & 0.387 & 0.144 & \textbf{0.991} \\
    \texttt{qwen3-0.6B}     & 0.641 & 0.349 & 0.227 & 0.902 \\
    \texttt{bge-m3}         & 0.604 & 0.288 & 0.237 & 0.813 \\
    \texttt{LaBSE}          & 0.511 & 0.314 & 0.295 & 0.779 \\
    \texttt{granite-278m}   & 0.529 & 0.289 & 0.318 & 0.757 \\
    \texttt{nomic-v2-moe}   & 0.580 & 0.239 & 0.215 & 0.727 \\
    \texttt{SapBERT}        & 0.352 & 0.142 & 0.424 & 0.494 \\
    \texttt{e5-large-instruct} & 0.237 & 0.021 & 0.481 & 0.239 \\
    \texttt{gte-base}       & 0.047 & 0.003 & 0.782 & 0.000 \\
    \bottomrule
  \end{tabular}
  \caption{Alias-graph benchmark leaderboard. ``pubR@10'' is per-publication
  Recall@10; ``conf'' is publication-lens confusion rate.
  \texttt{embeddinggemma} wins (MRS $0.991$, CI $[0.86,1.00]$).
  Source: alias-graph key findings, headline table / \texttt{headline\_numbers.csv}.}
  \label{tab:ag_board}
\end{table}

% CUT-NOTE (round 5): the two radar figures (cp_fig14, ag_fig10) were removed to
% meet the 8-page float budget; their only beat ("where each model wins") is
% carried by the leaderboard paragraph CUT-NOTE clause + Tables 1-2 + cp_fig17.

% ===========================================================================
\section{Analysis}
\label{sec:analysis}
% ===========================================================================

Why does cross-lingual retrieval collapse here, and what does the failure tell
us about the fix?

\paragraph{The modal failure is a same-language sibling.}
The two traps---language bias and chemical confusability---compound. Across the
$257$ confused concept-lens cases, the single most common failure is a
\emph{same-language sibling} out-ranking the gold ($114/257 = 44.4\%$); siblings
account for $79.4\%$ of all confusions and same-language winners for $55.6\%$
(Figure~\ref{fig:ag_joint}). The encoder is simultaneously biased toward the
query's own language and toward the nearest wrong compound. This joint mode is
the spine of the rest of the section, which we now read benchmark by benchmark.

\begin{figure}[t]
  \centering
  \includegraphics[width=\linewidth]{ag_fig12_joint_failure_modes.png}
  \caption{Joint failure modes among confused cases (alias-graph): the modal
  failure is a same-language sibling ($44.4\%$). Source: alias-graph
  \texttt{extra\_joint\_failure}.}
  \label{fig:ag_joint}
\end{figure}

\paragraph{Availability sets the stage; a residual encoder bias remains.}
On the \emph{alias-graph} benchmark, each language retrieves its own-language
gold far better ($0.63$--$0.82$) than foreign-language gold ($0.35$--$0.47$), and
because English and French source patents dominate, $42\%$ of an English query's
gold is reachable in-English versus only $8$--$10\%$ for German, Spanish, and
Chinese.\footnote{Availability shares and own-vs-foreign recall, alias-graph key
findings ``mechanism.''} Availability therefore shapes \emph{which} language's
gold is reachable---but it does \emph{not} explain away the same-language bias.
Regressing each language's home advantage on its gold-availability share gives a
\emph{negative} slope ($-0.57$, $n{=}5$ languages, descriptive): the language
with the \emph{least} reachable gold (Chinese, $8\%$) carries the \emph{largest}
home advantage ($+0.47$), so the mean $+0.32$ home advantage is a residual
encoder bias, not an availability artifact (Figure~\ref{fig:ag_avail}). On the
\emph{cross-lingual} benchmark the same bias shows as language collapse:
low-resource query languages over-fetch their own language up to $49\times$ the
corpus base rate, and same-language \emph{noise} out-ranks the gold on $60\%$ of
queries (Figure~\ref{fig:cp_collapse}).

\begin{figure}[t]
  \centering
  \includegraphics[width=\linewidth]{ag_fig11_availability_residual.png}
  \caption{Home advantage vs.\ gold-availability share per language
  (alias-graph). The slope is $-0.57$ ($n{=}5$, descriptive): home advantage
  does \emph{not} track availability, so it is residual encoder bias. Source:
  alias-graph \texttt{extra\_availability\_residual}.}
  \label{fig:ag_avail}
\end{figure}

\begin{figure}[t]
  \centering
  \includegraphics[width=\linewidth]{cp_fig09_10_collapse.png}
  \caption{Language collapse (cross-lingual benchmark). \emph{Left:}
  same-language over-representation by query language, up to $49\times$ the
  corpus base rate. \emph{Right:} the language of the documents that bury the
  gold; same-language noise out-ranks the gold on $60\%$ of queries. Source:
  chem-patents key findings, Figs.~9--10.}
  \label{fig:cp_collapse}
\end{figure}

\paragraph{Structure-style questions are the trap.}
\emph{How} a query is phrased decides whether the model falls into a confuser.
\emph{Structure}-style questions are the trap (Recall@10 $0.26$, confusion
$51\%$) because structure descriptions are exactly what sibling compounds
share, whereas \emph{role} questions are far safer (Recall@10 $0.60$, confusion
$25\%$); a language-independent \emph{formula token} (e.g.\ H$_2$S, CO$_2$)
measurably helps ($p<0.01$), even in Chinese (Figure~\ref{fig:ag_qtype}). The
trap is also where \emph{no} encoder helps: of the $16/132$ ($12\%$) queries
missed by all nine models, $14$ are structure questions about sibling compounds
(predominantly French and Chinese)---the universal-blind core needs
chemistry-aware help, not more encoders (\S\ref{sec:deployment}).

\begin{figure}[t]
  \centering
  \includegraphics[width=\linewidth]{ag_fig6_question_type_effect.png}
  \caption{Structure questions are the trap (R@10 $0.26$, confusion $51\%$) vs.\
  role questions ($0.60$, $25\%$). Source: alias-graph key findings, Fig.~6.}
  \label{fig:ag_qtype}
\end{figure}

\paragraph{Bias and inconsistency are related, but the robust signal is one.}
The more a model leans on its home language, the less consistent it is across
languages---$r(\text{home advantage},\text{RBO})=-0.85$ (cross-lingual) and
$r(\text{same-language share},\text{RBO})=-0.87$ (alias-graph), two related bias
proxies. We report these as \emph{descriptive} observations on $n{=}9$ models:
both relationships are driven by the two degenerate encoders and do \emph{not}
survive dropping them (the over-representation$\sim$CLIR correlation flips sign
and the home-advantage$\sim$RBO correlation collapses on $n{=}7$).\footnote{Drop-the-collapsers
check, chem-patents \texttt{extra\_correlation\_robustness}.} We therefore do not
treat them as mechanisms; the one correlation that \emph{is} robust is the
separability link below.

\paragraph{It is a separability deficit, so re-ranking cannot fix it.}
Behind the confusion, the cosine score barely separates gold from sibling:
$\mathrm{AUC}(\text{gold}>\text{look-alike})=0.55$ for confused queries versus
$0.70$ otherwise (Figure~\ref{fig:ag_sep}). On the cross-lingual benchmark the
same story holds across models, and this is the section's load-bearing
mechanism: $r(\text{cross-language AUC},\text{CLIR@10})=+0.96$
(Figure~\ref{fig:cp_sep}), and unlike the two bias correlations its \emph{sign}
is robust. Across model-level bootstrap resamples the separability--floor
correlation is positive in $99.97\%$ of draws (point $r=0.96$); with only seven
non-degenerate models the confidence interval is correspondingly wide
($[0.73, 1.00]$), so we report the sign as the robust finding rather than its
magnitude (Appendix~\ref{app:robust}). ``Robust'' here means \emph{sign-robust
under resampling}, not tightly estimated. Foreign gold is \emph{under-scored},
not merely mis-ordered.
This is the crux: a monolingual re-ranker reads a ranked list it cannot repair,
because the right foreign document was never scored high enough to appear---and
RRC bounds exactly how much it could repair. The budget curve (\S\ref{ssec:cp})
already flattens by the knee $K^{\!*}{=}5$, a top-100 re-ranker recovers at most
$\mathrm{RRC@}100=0.74$, and the floor $L_\infty=5.84\%$ of foreign twins is
unrecoverable by \emph{any} re-ranker. The ARI decomposition (\S\ref{ssec:cp})
shows this floor is the \emph{only} alignment-bound part: for the deployed model,
after a cheap top-100 re-rank, just $\mathrm{ARI@}100=0.229$ of the remaining gap
requires alignment (tied with \texttt{qwen3-0.6B} at $0.233$ for the lowest
re-ranker-irreducible residual), while \texttt{embeddinggemma} retains the
smallest alignment-only floor ($L_\infty=0.058$). The lever is at the embedding
level. We frame this descriptively rather than causally: cross-language AUC and
overall Recall@10 are strongly collinear across our seven models, and once
Recall@10 is partialled out the separability--CLIR association is no longer
significant (partial $r=+0.29$, $p=0.57$, n.s.), so separability is a descriptive
correlate of the cross-lingual floor, not an effect net of general retrieval
capability (Appendix~\ref{app:robust}).

\begin{figure}[t]
  \centering
  \includegraphics[width=\linewidth]{cp_fig11_separability.png}
  \caption{Per model, cross-language gold (red) is systematically harder to
  separate from non-gold than same-language gold (green); the model-level
  cross-language-AUC--CLIR@10 correlation is $+0.96$ ($+0.958$ on the $n{=}7$
  non-degenerate models). The \emph{sign} is robust under resampling
  ($P(r>0)=0.9997$); at $n{=}7$ the $95\%$ CI is wide ($[0.73, 1.00]$), so we
  report the sign, not the magnitude (Appendix~\ref{app:robust}). Source:
  chem-patents key findings, Fig.~11; robustness from
  \texttt{extra\_correlation\_robustness}.}
  \label{fig:cp_sep}
\end{figure}

\begin{figure}[t]
  \centering
  \includegraphics[width=\linewidth]{ag_fig8_confusion_is_separability.png}
  \caption{Confusion is a separability collapse: AUC $0.55$ confused vs.\ $0.70$
  otherwise. Source: alias-graph key findings, Fig.~8.}
  \label{fig:ag_sep}
\end{figure}

% ===========================================================================
\section{Deployment Recommendation}
\label{sec:deployment}
% ===========================================================================

The metrics yield a single, actionable decision for a search team.

\paragraph{Deploy \texttt{embeddinggemma}: the capability corner of the cost
frontier.}
We recommend \texttt{embeddinggemma} because it leads the capability axes
\emph{individually}---it tops cross-lingual Recall@10 ($0.50$) and separability
and is the best twin-finder (median first-foreign rank $5$)---and because it sits
on the cost-vs-capability frontier as its \emph{unique maximum-CLIR corner}: no
model is both cheaper-to-read \emph{and} higher-CLIR
(Figure~\ref{fig:cost_frontier}). It is \emph{not} the cheapest deployable model:
\texttt{bge-m3} reads shallower (XRC50 $2.0\times$ vs.\ $3.5\times$) at a
still-admitted CLIR@10, so the decision is a \emph{frontier choice}---pick the
capability corner unless the reading budget is binding, in which case
\texttt{bge-m3} is the cheaper-to-read frontier alternative over the stated
$\tau$-band ($\tau \in [0.33, 0.44]$; below $\tau{\approx}0.33$ the cheaper-reader
pick flips to \texttt{granite-278m}, while only the \texttt{embeddinggemma}
max-CLIR corner is $\tau$-invariant; \S\ref{ssec:cp}). It also carries the
smallest alignment-only floor ($L_\infty=0.058$) and is tied with
\texttt{qwen3-0.6B} for the lowest $\mathrm{ARI@}100$ ($0.229$ vs.\ $0.233$, gap
$0.004$ with a $95\%$ CI straddling zero; \S\ref{ssec:cp}), so the part of its
cross-lingual gap a re-ranker cannot close is among the smallest of the
deployable field. We deliberately do
\emph{not} rest the recommendation on the composite: \texttt{embeddinggemma}'s
rank ranges $[1,4]$ across aggregation schemes (\S\ref{ssec:boards}), so a
composite-based argument would be fragile. Per-axis dominance and the frontier
are not.

\paragraph{The capability corner is route-dependent: a per-route router is
upside headroom.}
Recomputing the $(\mathrm{XRC50}_\ell,\mathrm{CLIR@10}_\ell)$ frontier
\emph{per query language} $\ell$ shows the capability corner is not stationary:
across the five routes there are \emph{three} distinct max-CLIR corners
(en$\rightarrow$\texttt{qwen3-0.6B}, de/es/zh$\rightarrow$\texttt{embeddinggemma},
fr$\rightarrow$\texttt{nomic-v2-moe}), and a recall-only dashboard's pick
\emph{flips} on $2/5$ routes (en, fr; Figure~\ref{fig:per_route}). This nuances,
but does \emph{not} overturn, the single-model recommendation. First,
\texttt{embeddinggemma} remains the right default: it is the global capability
corner (max pooled CLIR@10 $=0.50$), owns the smallest alignment-only floor
($L_\infty=0.058$) and ties \texttt{qwen3-0.6B} for the lowest $\mathrm{ARI@}100$
($0.229$ vs.\ $0.233$, gap CI straddles zero), and per-route wins
$3/5$ routes (de, es, zh)---including the two hardest cross-only routes (Spanish
has \emph{zero} same-language gold; Chinese is the thinnest)---while the
recall-only dashboard picks it on all five. Second, a per-route router that swaps
in \texttt{qwen3-0.6B} for en and \texttt{nomic-v2-moe} for fr is a plausible
\emph{upside}, consistent with the oracle/ensemble headroom below, but on our
per-language sample sizes we report it as headroom, not a recommendation. The
route corners rest on thin per-language samples (de $n_{\text{same}}{=}7$, zh
$n_{\text{same}}{=}2$, es $n_{\text{same}}{=}0$; cross-side $n{=}22$--$34$): the
defensible spine is the robust per-route $\mathrm{CLIR@10}_\ell$ axis, the
frontier membership, and the \emph{existence} of corner movement, whereas the
per-route XRC axis is \emph{indicative} only and es is XRC-undefined (zero
same-language gold) and never imputed. We therefore present routing as the honest
upside of the frontier, not a reversal of the deploy-one-model decision.

\begin{figure}[t]
  \centering
  \includegraphics[width=\linewidth]{cp_fig23_per_route_frontier.png}
  \caption{Per-route cost frontier: the $(\mathrm{XRC50}_\ell,
  \mathrm{CLIR@10}_\ell)$ frontier recomputed per query language. The capability
  corner moves---three distinct max-CLIR corners across five routes
  (en$\rightarrow$\texttt{qwen3-0.6B}, de/es/zh$\rightarrow$\texttt{embeddinggemma},
  fr$\rightarrow$\texttt{nomic-v2-moe})---and the recall-only pick flips on $2/5$
  routes (en, fr). The CLIR@10$_\ell$ axis is robust; the per-route XRC axis is
  \emph{indicative} (es undefined, $n_{\text{same}}{=}0$, not imputed).
  \texttt{embeddinggemma} stays the single-model default (wins $3/5$ routes incl.\
  the hardest); a per-route router is upside headroom. Source: chem-patents
  \texttt{extra\_per\_route\_frontier}.}
  \label{fig:per_route}
\end{figure}

\paragraph{Report robustness next to recall.}
We recommend reporting XRC, RRC, CLIR@10, and language-parity \emph{alongside}
recall, never average recall alone. The pure ``no-home'' Spanish queries are the
reason: a dashboard that averages over them will look healthy while the deployed
system silently fails on cross-jurisdiction search.

\paragraph{Do not reflexively ensemble.}
There \emph{is} real headroom---the oracle (any of nine models) reaches CLIR@10
$0.61$ on the cross-lingual benchmark (Figure~\ref{fig:cp_ens}) and $88\%$
versus the best single model's $76\%$ on the alias-graph benchmark, with Chinese
having the largest recoverable headroom (Figure~\ref{fig:ag_ens}). But this
headroom is \emph{not free}: untuned reciprocal-rank fusion \citep{rrf2009} did
\emph{not} beat the best single model. The two facts coexist---naive fusion
loses, yet a score-aware learned combiner or per-language routing (for the
homeless Spanish and the low-resource Chinese) could win. And more encoders will
not close the gap on the hardest queries: the $16/132$ ($12\%$) universal-blind
core is overwhelmingly structure questions about sibling compounds ($14/16$;
\S\ref{sec:analysis}), which need chemistry-aware help, not another retriever.

\begin{figure}[t]
  \centering
  \includegraphics[width=\linewidth]{cp_fig12_ensemble_headroom.png}
  \caption{Best-single vs.\ fusion vs.\ router vs.\ oracle (cross-lingual);
  oracle CLIR@10 $0.61$, but untuned RRF does not beat the best single model.
  Source: chem-patents key findings, Fig.~12.}
  \label{fig:cp_ens}
\end{figure}

\begin{figure}[t]
  \centering
  \includegraphics[width=\linewidth]{ag_fig7_ensemble_headroom.png}
  \caption{Oracle vs.\ best single model, per language (alias-graph); $88\%$ vs.\
  $76\%$, Chinese has the largest headroom.
  Source: alias-graph key findings, Fig.~7.}
  \label{fig:ag_ens}
\end{figure}

\paragraph{Budget the re-ranker by the knee.}
If a re-ranker is used, the RRC budget curve (Figure~\ref{fig:rrc_budget}) sets
its depth: most recoverable payoff sits in a shallow top-$K$---the knee is
$K^{\!*}{\approx}5$ for the deployed \texttt{embeddinggemma} and $\le 20$ for
nearly every non-degenerate model---so a deeper candidate pool past the knee buys
almost nothing. There is no benefit in paying to re-rank a top-1000 pool when the
returns flatten by the top-20.

\paragraph{Align, do not re-rank.}
Because foreign gold is under-scored at the embedding level
(\S\ref{sec:analysis}), the durable fix is representation alignment at indexing
time. A monolingual re-ranker at query time operates on a candidate list that
already excludes the right foreign document---and RRC says how much: for the
deployed model, $L_\infty = 1-\mathrm{RRC@1000}=5.84\%$ of foreign twins are
unrecoverable by \emph{any} re-ranker (the alignment-only floor), and a realistic
top-100 re-ranker leaves $25\%$ on the table.

\paragraph{Budget rule.}
Finally, the MT-of-question null result ($-0.044$, $p=0.13$) tells us where to
spend a finite human-translation budget: \emph{machine-translate the query, but
human-translate the corpus}. The expensive resource buys the most where it
matters---the documents being searched, not the questions. This re-derives the
classic query-vs-document translation rule \citep{oard1998,salehpecina2020} for
embedding retrieval over patents, where our null quantifies the query-side MT
cost as insignificant.

% ===========================================================================
\section*{Limitations}
% ===========================================================================

\paragraph{Scale and statistical reach.}
We consolidate the resample interval or sign-stability vote of every
load-bearing small-$n$ scalar---the separability $r$, the three frontier XRC50
medians, the \texttt{embeddinggemma}-vs-\texttt{qwen3-0.6B} ARI@100 tie, and the
separability partial-$r$---in Appendix~\ref{app:robust}, so the honest intervals
are reported in one place rather than buried.
The benchmarks comprise 132 (alias-graph) and 137 (cross-lingual) questions over
24 compounds; some language-pair cells are thin, so the directional CLIR matrix
should be read as indicative rather than definitive in the sparsest cells. For
the same reason we report XRC50 (the median) as the robust headline and treat
the deeper $D_{90}/D_{95}$ coverages as right-censored lower bounds, not point
estimates; and the negative availability slope ($-0.57$) is a \emph{descriptive}
regression on five language points, not an inferential test. Two further
cross-model correlations are reported as descriptive/motivating only, not
established effects: the two-tax non-redundancy
($\rho{=}{-}0.59$, $n{=}7$, $p{=}0.16$) and the cheapest-reader trap
($\rho{=}{+}0.29$, $n{=}7$, $p{=}0.53$) are both \emph{non-significant} on seven
non-degenerate models. The RRC budget objects (knee $K^{\!*}$, floor $L_\infty$)
and the ARI decomposition (which closes to $1.0$ for all nine models), by
contrast, are regression-checked and stated firmly.

\paragraph{Per-route frontier thinness.}
The per-route capability corners (Figure~\ref{fig:per_route}) are estimated on
thin per-language samples (de $n_{\text{same}}{=}7$, zh $n_{\text{same}}{=}2$, es
$n_{\text{same}}{=}0$; cross-side $n{=}22$--$34$). The per-route XRC axis is
therefore \emph{indicative} only, and the Spanish route's XRC is undefined (no
same-language gold) and is never imputed. Accordingly we report the per-route
router as \emph{headroom}, not a recommendation: the deployed decision is a single
model, and the robust per-route signals are the CLIR@10$_\ell$ axis, the frontier
membership, and the existence of corner movement.

\paragraph{Domain transfer.}
All results are chemistry-patent-specific. We frame that specificity as the
contribution---a real deployment niche underserved by general benchmarks---but
we do not yet provide a general-domain MMTEB/MIRACL/NeuCLIR companion
measurement; it is backlogged.
% TODO trace: no MMTEB / MIRACL / NeuCLIR transfer number computed yet; flagged
% in needs_eval as a candidate.

\paragraph{Composite validation.}
CLIR-MRS/MRS is a reporting convenience, not externally validated against a
downstream utility; the per-axis dominance is what the recommendation stands on.
A validation against a held-out human or downstream-RAG criterion is backlogged.

\paragraph{Judge dependence and gold equivalence.}
The generated queries and the auto-grader rely on an LLM judge. We provide human
validation of a sample (Appendix), but claim-level patent equivalence is not
lay-annotatable; a larger expert re-annotation, and a spot-check that the
parallel human-translated golds are claim-level equivalent, remain future work.

\paragraph{Severity law.}
Confusion severity is reported as a two-level sibling/parent split; the graded
ChEBI hop-distance decay law is future work, since the on-disk relation field is
binary.

\paragraph{Language coverage and a forthcoming causal probe.}
We cover five languages (en/de/es/fr/zh) with a single non-Latin script
(Chinese); low-resource scripts are out of scope here. Calibration-style
cross-lingual instruments and a conformal (guaranteed-coverage) version of XRC
are likewise left as future work; we do not claim them as results. Our
``align, not re-rank'' conclusion is correlational; the natural next step---which
XRC and RRC were built to measure---is to fit a per-language alignment map on one
model, re-retrieve, and recompute XRC50/RRC@100/ARI before and after. The ARI
decomposition (Figure~\ref{fig:ari}) hands this probe a precise before/after
target: the per-model $L_\infty$ floor / $\mathrm{ARI@}100$ is exactly the
quantity an alignment intervention should move while re-ranking cannot. This is
upside only; none of the results above depend on it, and we flag it as
forthcoming.

% ===========================================================================
\section{Conclusion}
% ===========================================================================

For a chemistry-patent search team deploying one multilingual embedding model,
average Recall@10 is the wrong dashboard: it is inflated by same-language hits
and hides a cross-lingual collapse (best CLIR@10 $0.50$; home advantage up to
$+0.55$; the best cross-lingual RBO any model reaches is $0.39$ on the
alias-graph benchmark and $0.19$ on the cross-lingual benchmark; a confusable
compound beating the gold on $14$--$78\%$ of queries). That collapse \emph{costs}:
finding a foreign twin takes scanning $\sim$3.5$\times$ the reading budget (XRC),
and a re-ranker cannot fully recover it---the budget curve knees at $K^{\!*}{=}5$
and leaves a floor $L_\infty=0.058$ that is the only part of the gap a re-ranker
cannot move, recoverable by alignment alone (RRC). Our two
content-controlled, patent-grounded benchmarks and our CLIR robustness suite make
the collapse measurable, place the deployable models on a cost-vs-capability
frontier on which \texttt{embeddinggemma} is the Pareto-optimal capability corner
(\texttt{bge-m3} the cheaper-to-read alternative), and trace the cause to an
embedding-level separability deficit
($r(\text{cross-language AUC},\text{CLIR@10})=+0.96$, robust)---so the fix is
alignment at index time, not re-ranking at query time, and the budget rule is to
machine-translate the query and human-translate the corpus. Reporting robustness
next to recall, and treating cross-lingual retrieval as a representation problem,
are the two habits we recommend any multilingual patent-search team adopt.

\bibliography{custom}

\appendix
\section{Appendix}

\paragraph{Robustness of the load-bearing small-$n$ scalars.}
\label{app:robust}
Table~\ref{tab:robust} consolidates the resample interval (or sign-stability
vote) of every small-$n$ scalar the body leans on. Three points are deliberate.
(i) The separability $r$ is reported by its \emph{sign} ($P(r>0)=0.9997$), not
its magnitude: at $n{=}7$ the $95\%$ CI is wide. (ii) The three frontier XRC50
medians have \emph{finite} but wide CIs (censored-draw fraction $=0$;
median-of-discrete-depth bootstrap), so we report the depth \emph{ordering}, not
CI tightness. (iii) The \texttt{embeddinggemma}-vs-\texttt{qwen3-0.6B} ARI@100
gap CI straddles zero (a tie), and the separability--CLIR association is no longer
significant once Recall@10 is partialled out---so separability is a descriptive
correlate, not an effect net of capability. Source:
\texttt{extra\_robustness\_appendix/\{robustness\_table.csv, summary.json\}}.

\begin{table}[t]
  \centering
  \small
  \begin{tabular}{lrl}
    \toprule
    scalar & point & resample / stability ($n$) \\
    \midrule
    separability $r$ & $0.96$ & $P(r{>}0){=}0.9997$; CI $[0.73,1.00]$ ($n{=}7$) \\
    XRC50 \texttt{egemma} & $3.5$ & CI $[0.91,12.0]$ (finite, wide) \\
    XRC50 \texttt{bge-m3} & $2.0$ & CI $[0.53,7.0]$ (finite, wide) \\
    XRC50 \texttt{granite} & $1.25$ & CI $[0.28,12.25]$ (finite, wide) \\
    ARI@100 gap (q$-$e) & $0.004$ & CI $[-0.174,0.176]$; $P{=}0.519$ \\
    partial $r\,|\,$R@10 & $0.29$ & $p{=}0.57$ (n.s.); zero-order $0.96$ \\
    \bottomrule
  \end{tabular}
  \caption{Resample robustness of the body's small-$n$ scalars. Separability $r$
  is reported by its sign (wide CI at $n{=}7$); the three frontier XRC50 CIs are
  finite but wide (censored-draw fraction $=0$); the ARI@100 gap
  (\texttt{qwen3-0.6B}$-$\texttt{embeddinggemma}) CI straddles zero (a tie); and
  the separability--CLIR partial $r$ controlling Recall@10 is non-significant.
  Source: chem-patents \texttt{extra\_robustness\_appendix}.}
  \label{tab:robust}
\end{table}

\paragraph{Reproducibility.}
The full ten-round derivations behind the headline numbers, the per-language and
per-lens tables, and all figures are produced by the scripts under
\texttt{reports/runs/\{chem\_patents,alias\_graph\}/experimental\_codes/} and
summarised in the two \texttt{key\_findings/EXECUTIVE\_SUMMARY.md} files this
paper cites.

\paragraph{Generation-pipeline validation (C5).}
A human annotator reviewed a sample of generated queries; the quantitative
human-rating summary and the LLM auto-grader's strictness relative to humans are
summarized in the system description and will be folded in here once dumped to
\texttt{reports/}.
% TODO trace: human-eval numbers (mean rating, items reviewed, rejections,
% auto-grader strictness) currently in source slides only, not under reports/.

\end{document}
